% Pursuing Stacks (cotes 134-2 to 134-8): archive cote 134-7, site page 66 (PDF page 67), typescript page 508.
% typescript. Draft compared: ps6.tex:3218-3296.
% AI-transcribed by Claude Opus 5.5 (claude-opus-5-5, adaptive thinking, effort high), Message Batches API, batch ps4 (msgbatch_012Aoc1NuskATPsyBeKSYK7J), 2026-11-12; not yet checked by eye.
\scanpage{134-7/66}
\mspage{508}

121.    The tentative approach towards\xout{a} defining and studying "schematic"
homotopy types I have been following lately relies heavily on a suitable
notion of "linearization" of such homotopy types.  One \xout{may} can imagine
that \deleted{very much}\inserted{many} different appr\xout{w}oaches (for instance via De Rham complexes
with devided powers,  or additive small categories with diagonal maps)
\margin{$\times$}may be devised for "schematic" homotopy types,  but in \uline{any case} it seems
likely that a suitable notion of linearization will play an important
role\xout{,}.  \xout{andx}It may be worthwhile therefore to try and pin down the wished-%
for main features of such a theory,  with the hope maybe of getting an
axiomatic description for it,  with a corresponding unicity statement.
Before doing so,  the first thing to do seems to review some main formal
features of linearization for ordinary ("discrete") homotopy types.

Recall the definition

(1) \qquad Hotab $\overset{\mathrm{def}}{=}$ D.(Ab) = derived category of the category of
\xout{chain} abelian chain complexes\hand{,} with
respect to quasi-isomorphisms\hand{,}

and the two canonical functors

(2) \qquad \diagram{\begin{tikzcd} \mathrm{Hot} \ar[r, shift left, "\hand{LH.}"] & \mathrm{Hotab} \ar[l, shift left, "\deleted{\illegible{1}}\hand{K_\pi}"] \end{tikzcd}} \quad ,

Where \deleted{\illegible{1}.}\inserted{LH.} is the "abelianization functor",  and \deleted{DP}\inserted{$K_\pi$} is\inserted{\deleted{$K_\pi$}}\inserted[typed]{defined via} the Kan-Dold-Puppe
functor,  associating to a chain complex the \xout{homotopy type} corresponding
semisimplicial abelian group .  The diagram (2) may be viewed\inserted{(up to equivalence)} as deduced
from the corresponding diagram

(2') \qquad \diagram{\begin{tikzcd} \Delta^{\wedge} \ar[r, shift left, "W"] & \Delta^{\wedge}_{ab} \ar[l, shift left, "DP"] \end{tikzcd}} \quad ($\cong$ Ch.(Ab))

by passing to the suitable localized categories Hot and Hotab.  In the
diagram (2')  \xout{DPxisxright}\inserted[typed]{W is left} adjoint to DP ,  which is now just the forget-%
ful functor.   One main fact about (2) is

(3) \qquad \over{K}{$K_\pi$} is right adjoint to LH. \quad ,

which is just a neater way \over{to}{for} express\inserted{ing} the familiar fact that for given
abelian group $\pi$ and  natural integer n,  the object K(\hand{$\pi$},n) in  H\over{o}{\unsure{\v{o}}}t
(namely the image by \deleted{K}\inserted{$K_\pi$} of the chain complex\inserted{$\pi[n]$} reduced to \hand{$\pi$} in degree n)
represents the cohomology functor
\[ X \mapsto H^n(X, \hand{\pi}) \simeq \mathrm{Hom}_{\mathrm{Hotab}}(LH.(X), \pi[n]) \quad . \]
The functor \deleted{K}\inserted{$K_\pi$} may be called the "\uline{Eilenberg-Maclane functor}",  as it's
values are immediate generalizations of the Eilenberg-Maclane objects
\deleted{\unsure{K($\pi$,n)}}.

\margin{66}
